\chapter*{Tóm Tắt}
\addcontentsline{toc}{chapter}{Tóm Tắt}

Khóa luận đề xuất và xây dựng \system{} -- một hệ thống luyện phỏng vấn IT cá nhân hóa tích hợp \textbf{Knowledge Tracing (KT)} cho ba vai trò dữ liệu phổ biến: Data Analyst (DA), Data Scientist (DS) và Data Engineer (DE). Hệ thống giải quyết hai hạn chế của các nền tảng luyện phỏng vấn hiện nay: thiếu cá nhân hóa theo năng lực và không theo dõi được quỹ đạo kiến thức của người học.

Về dữ liệu, khóa luận xây dựng question bank đa dạng (5 loại câu hỏi, 17 Knowledge Components) bằng pipeline kết hợp thu thập từ các nguồn công khai và sinh tự động bằng LLM kèm kiểm định chất lượng; đồng thời thiết kế pipeline mô phỏng tương tác dựa trên Item Response Theory (IRT) để tạo bộ dữ liệu huấn luyện KT (1.000 người dùng ảo, $\sim$40.000 tương tác). Về mô hình, khóa luận huấn luyện và so sánh ba họ KT -- BKT, DKT, SAKT -- cùng biến thể \textit{quality label} liên tục. Kết quả trên tập kiểm tra cho thấy mọi mô hình KT vượt baseline (ROC-AUC $\approx$ 0.68 so với 0.50); biến thể \textbf{DKT-Quality} đạt PR-AUC và RMSE tốt nhất trong nhóm neural.

Về hệ thống, khóa luận thiết kế cơ chế cập nhật năng lực thời gian thực bằng cách kết hợp dự đoán KT (70\%) và trung bình trượt mũ EMA (30\%), tích hợp vào engine gợi ý theo chiến lược \textit{zone of proximal development} và bài kiểm tra thích nghi hai giai đoạn (Multi-Stage Testing). Toàn bộ được hiện thực thành ứng dụng web hoàn chỉnh (React + FastAPI + Streamlit) với chấm điểm tự động bằng Gemini LLM và cơ chế chấm dự phòng cục bộ.

\textbf{Từ khóa:} Knowledge Tracing, Deep Knowledge Tracing, Item Response Theory, Hệ thống gợi ý, Computerized Adaptive Testing, Luyện phỏng vấn IT.

\vspace{0.6cm}
\begin{center}\rule{0.5\textwidth}{0.4pt}\end{center}
\vspace{0.3cm}

\noindent\textbf{\large Abstract}

\smallskip
This thesis proposes and implements \system{}, a personalized IT interview-practice platform powered by \textbf{Knowledge Tracing (KT)} for three popular data roles: Data Analyst, Data Scientist, and Data Engineer. It addresses two limitations of existing platforms: the lack of competency-based personalization and the inability to track a learner's knowledge trajectory. A diverse question bank (5 question types, 17 Knowledge Components) is built through a hybrid pipeline of public-source collection and LLM generation with quality control, while an IRT-based simulation pipeline produces the KT training data (1{,}000 virtual users, $\sim$40{,}000 interactions). Three KT model families -- BKT, DKT, SAKT -- are trained and compared, including a continuous \textit{quality-label} variant; all KT models beat the baseline (ROC-AUC $\approx$ 0.68 vs.\ 0.50), with DKT-Quality achieving the best PR-AUC and RMSE among neural models. A real-time skill-update mechanism blends KT prediction (70\%) with an exponential moving average (30\%), feeding a zone-of-proximal-development recommender and a two-stage adaptive test. The whole pipeline is realized as a complete web application (React + FastAPI + Streamlit) with automatic LLM-based grading and a local fallback grader.

\noindent\textbf{Keywords:} Knowledge Tracing, Deep Knowledge Tracing, Item Response Theory, Recommendation System, Computerized Adaptive Testing, IT Interview Preparation.
